% Type System: formal grammar of theorem statements in the Open Beilinson tower.

\providecommand{\BeilLevel}{\mathrm{Lvl}}
\providecommand{\Quad}{\mathrm{Q}}
\providecommand{\Pres}{\mathrm{P}}
\providecommand{\Hyp}{\mathrm{H}}
\providecommand{\TypeSig}{\mathfrak{T}}
\providecommand{\Open}{\mathsf{Open}}
\providecommand{\KSDual}{\mathsf{KSDual}}
\providecommand{\Catl}{\mathsf{Cat}}
\providecommand{\Chnl}{\mathsf{Chain}}
\providecommand{\HypAug}{\mathsf{Aug}}
\providecommand{\HypComp}{\mathsf{Comp}}
\providecommand{\HypFinDimGr}{\mathsf{FDG}}
\providecommand{\HypKosz}{\mathsf{Kosz}}
\providecommand{\HypAmb}{\mathsf{Amb}}
\providecommand{\HypGen}{\mathsf{Gen}}
\providecommand{\HypClass}{\mathsf{Cl}}
\providecommand{\HypProConil}{\mathsf{ProConil}}
\providecommand{\HypJad}{\mathsf{J\text{-}adic}}
\providecommand{\HypWtComp}{\mathsf{WtComp}}
\providecommand{\HypCoder}{\mathsf{Coder}}
\providecommand{\HypMC}{\mathsf{MC}}
\providecommand{\HypMod}{\mathsf{Mod}}
\providecommand{\HypClut}{\mathsf{Clut}}
\providecommand{\HypHall}{\mathsf{Hall}}
\providecommand{\HypIntegral}{\mathsf{Int}}
\providecommand{\HypDescent}{\mathsf{Desc}}
\providecommand{\Recon}{\mathrm{Rec}}
\providecommand{\meet}{\sqcap}
\providecommand{\joinp}{\sqcup}
\providecommand{\Ab}{\mathcal{A}_b}

\chapter[Type system]{Type system: the grammar of theorem statements}
\label{app:type-system}

The Open quadrant of the four-quadrant theory organises chiral
algebra on $(X,D,\tau)$ as a five-step Beilinson tower with a
reconstruction theorem at each interface. Every theorem in this
volume carries a four-tuple
\[
 \TypeSig
 = (\Quad,\Pres,\BeilLevel,\Hyp),
\]
the \emph{type signature}: quadrant, presentation, level, hypothesis
package. Every cross-level equality invokes a named reconstruction
theorem with explicit hypothesis package; suppressing the package
collapses the quadrant grid (shadow-object distinction) and turns a
forgetful image into a false equivalence.

This appendix fixes the grammar. Section~\ref{sec:type-levels} fixes
the level alphabet and gives an inscription example at each level.
Section~\ref{sec:type-quadrants} fixes the quadrant alphabet and the
$\mathbb{Z}/2$ involution defining $\KSDual$.
Section~\ref{sec:type-presentations} fixes the chain-level /
$(\infty,1)$-categorical dichotomy as equally load-bearing
presentations. Section~\ref{sec:type-hypothesis} catalogues the
hypothesis-package primitives and their containment partial order.
Section~\ref{sec:type-composition} gives the composition rule for
packages along reconstruction theorems and the meet/join.
Section~\ref{sec:type-lattice} draws the Hypothesis-Package Lattice
as a TikZ figure on the reconstruction-theorem nodes.

\section{Levels}
\label{sec:type-levels}

\begin{definition}[Beilinson levels]
\label{def:type-levels}
The Open Beilinson tower is the linearly ordered five-element set
\[
 \BeilLevel
 \in\{0,1,2,3,4,5\},
\]
indexing the carriers
\[
\begin{array}{c|l|l}
\BeilLevel & \text{name} & \text{carrier} \\
\hline
0 & \text{primitive} & \mathcal C^{\mathrm{op}}\text{ on }(X,D,\tau)\text{: factorisation dg-cat} \\
1 & \text{chart} & \Ab=\mathrm{End}_{\mathcal C}(b),\ b\in\mathcal C^{\mathrm{op}} \\
2 & \text{twisting} & B(\Ab)=T^c(s^{-1}\bar\Ab) \\
3 & \text{closed sector} & Z^{\mathrm{der}}_{\mathrm{ch}}(\Ab)\simeq\ChirHoch^\bullet(\Ab) \\
4 & \text{operator} & \mathsf{Line}_{\Ab},\ \mathsf{Brane}_{\Ab} \\
5 & \text{scalar} & (\kappa\text{-tuple},Z_g,\Phi_N).
\end{array}
\]
The level alphabet is exhaustive on the Open quadrant: every Vol~I
theorem inscribes its conclusion at exactly one level.
\end{definition}

The level signature is recoverable from the carrier by inspection;
a theorem whose conclusion compares carriers at different levels is
called \emph{cross-level}, and the equality is then an image
equality under a reconstruction theorem.

\begin{example}[Level inscriptions]
\label{ex:type-level-inscriptions}
\begin{enumerate}[label=\textup{(L\arabic*)}]
\item \emph{Level 0.} The factorisation dg-cat
$\mathcal C^{\mathrm{op}}$ on $(X,D,\tau)$ is itself the inscription;
the Morita equivalence $\mathcal C^{\mathrm{op}}\simeq
\mathrm{Mod}_{\Ab}^{\mathrm{fact}}$ at chosen vacuum~$b$ is a
\emph{cross-level} statement $0\leftrightarrow 1$.
\item \emph{Level 1.} Theorem~A is the level-$1\leftrightarrow 2$
reflection $K^2\simeq\mathrm{id}$ on $\mathrm{Kosz}(X)$.
\item \emph{Level 2.} The bar coalgebra $B(\Ab)$ is the carrier of
the universal twisting morphism
$\pi_{\Ab}\colon B(\Ab)\to\Ab$ underlying Theorem~B inversion.
\item \emph{Level 3.} Theorem~H states
$\ChirHoch^\bullet(\Ab)\subset\{0,1,2\}$ on the Koszul locus at
generic level.
\item \emph{Level 4.} Theorem~D states
\(\mathrm{obs}_g(\Ab)=\kappa(\Ab)\lambda_g\) as an
$H^\bullet(\overline{M}_{g,n})$ class.
\item \emph{Level 5.} Theorem~C states the family-stratum ceiling
$K^\kappa(\Ab)=\kappa(\Ab)+\kappa(\Ab^!)$ on the
$\kappa$-stratification matrix.
\end{enumerate}
\end{example}

\section{Quadrants}
\label{sec:type-quadrants}

\begin{definition}[Quadrant alphabet]
\label{def:type-quadrants}
The four-quadrant grid
\[
 \Quad
 \in\bigl\{\Open,\ \mathrm{CY}\bigr\}
 \times\bigl\{\Catl,\ \Chnl\bigr\}
\]
factors the theory: $\Open$ is Vol~I (algebraic-curve
factorisation), $\mathrm{CY}$ is Vol~III (Calabi--Yau-category two-stage
$\Phi_d^{(\Sigma_{d-1},C)}=\mathrm{Sp}^{\mathrm{ch}}_{\Sigma_{d-1},C}\circ\Phi_d^{\mathrm{FA}}$).
The presentation factor $\{\Catl,\Chnl\}$ is treated in
\S\ref{sec:type-presentations}.
\end{definition}

\begin{definition}[Chiral involution and $\KSDual$]
\label{def:type-ksdual}
Let $\sigma\colon\Ab\mapsto\Ab^!$ denote the chiral / antichiral
Verdier--Koszul involution: $\sigma$ is the composition
$\Ab\stackrel{B}{\longrightarrow}B(\Ab)\stackrel{\mathbb D_{\Ran}}{\longrightarrow}
B(\Ab)^\vee\stackrel{\Omega}{\longrightarrow}\Ab^!$ on the Koszul
locus. The fixed sublocus
\[
 \KSDual
 :=\{\Ab\mid\Ab\simeq\Ab^!\}
\]
is the $\mathbb Z/2$-fixed locus on which the four-quadrant grid
degenerates into self-dual form. On $\KSDual$:
\begin{itemize}
\item Theorem~A holds as an equivalence (not adjunction with
comparison map);
\item Theorem~H is exact in $\{0,1,2\}$;
\item the five-archetype dichotomy stabilises;
\item the Universal Trace Identity
$h_{\Ab}=h_{\mathrm{LV}}/\mathcal N(\Ab)$, with level-$2$
Verdier-Ran-on-bar norm
$\mathcal N(\Ab)=\kappa(\Ab)+\kappa^!_{\mathrm{LV}}(\Ab)$
(distinct from the algebra-level $K^\kappa$), is the equivariant
signature.
\end{itemize}
\end{definition}

\section{Presentations}
\label{sec:type-presentations}

\begin{definition}[Presentation tag]
\label{def:type-presentation}
The presentation tag $\Pres\in\{\Catl,\Chnl\}$ records the lane in
which a theorem is stated.
\begin{itemize}
\item $\Pres=\Catl$ ($(\infty,1)$-categorical): derived stable
$\infty$-categories, Lurie HA, factorisation $\infty$-categories,
$\infty$-operads (Costello--Gwilliam, Francis--Gaitsgory). Proof
witnesses are $(\infty,1)$-functors, adjunctions, colimits,
monoidal structures.
\item $\Pres=\Chnl$ (chain-level): explicit complexes, named
differentials, witnessed homotopies, $L_\infty$-structures,
Mittag--Leffler towers, $J$-adic / pro-conilpotent / weight-completed
ambients. Proof witnesses are chain homotopies, explicit MC
elements, OPE poles.
\end{itemize}
Both presentations are load-bearing. A theorem inscribed in $\Catl$
is not the shadow of its $\Chnl$ avatar, and conversely; both are the
theorem viewed through different lenses (shadow-object distinction
disallows demoting either to the other).
\end{definition}

When a theorem requires both lanes (e.g.\ Theorem~B inversion off
$\KSDual$), the type signature carries
$\Pres=\Catl\meet\Chnl$ with the chain-level lane
ambient-qualified in $\Hyp$.

\section{Hypothesis-package primitives}
\label{sec:type-hypothesis}

\begin{definition}[Primitive hypotheses]
\label{def:type-hypotheses}
The hypothesis alphabet is a finite set of primitives, each a
condition on $\Ab$, on the bar coalgebra $B(\Ab)$, on the ambient
category, or on the stratum:
\[
\begin{array}{l|l}
\textbf{Primitive} & \textbf{Condition} \\
\hline
\HypAug & \Ab\in\Alg^{\fact,\aug}_X;\ \bar\Ab=\ker\epsilon. \\
\HypComp & \Ab\simeq\varprojlim_n\Ab/\bar\Ab^{\,n}\text{ (weak curvature).} \\
\HypFinDimGr & \dim_k(B(\Ab))_w<\infty\ (\forall w\in\mathbb Z_{\geq 0}). \\
\HypKosz & \mathrm{Ext}^{i,j}_{\mathcal C^{\mathrm{op}}}(b,b)=0\ (i\neq j); \\
       & \text{i.e. }\Ab\in\mathrm{Kosz}(X). \\
\HypAmb(\square) & \text{ambient }\square\in\{\Ch(\Vect),\HypCoder,\HypWtComp,\Pro,\HypJad\}. \\
\HypGen & \text{generic level: }k\notin\{-h^\vee\}\cup\{\text{admissible roots}\}. \\
\HypClass(\square) & \text{archetype label }\square\in\{\mathsf G,\mathsf L,\mathsf C,\mathsf M,\mathsf B\}. \\
\HypMC & \text{Maurer--Cartan element }\Theta_{\Ab}\text{ exists and converges.} \\
\HypMod & \text{modular trace data on the open category.} \\
\HypClut & \text{clutching data on }\overline{M}_{g,n}. \\
\HypHall & \text{Hall pairing on a half-density category.} \\
\HypIntegral & \text{integral form of the deformed enveloping object.} \\
\HypDescent & \text{descent data along a fibration of curve geometries.}
\end{array}
\]
A \emph{hypothesis package} is a finite conjunction
$\Hyp=\bigwedge_i \Hyp_i$ of primitives.
\end{definition}

\begin{example}[Theorem-A package]
\label{ex:type-thmA-package}
The Theorem~A package (\S\ref{sec:ainf2-koszul-reflection}, hypotheses
H1--H3) is
\[
 \Hyp_{\mathrm{A}}
 = \HypAug\wedge\HypComp\wedge\HypFinDimGr,
\]
with chain-level class~$\mathsf M$ requiring the strengthening
\[
 \Hyp_{\mathrm{A},\mathsf M}
 = \Hyp_{\mathrm{A}}\wedge\HypAmb(\HypWtComp\joinp\Pro\joinp\HypJad).
\]
\end{example}

\begin{example}[Theorem-H package]
\label{ex:type-thmH-package}
The Theorem~H package on the smooth Koszul locus is
$\Hyp_{\mathrm{H}}=\Hyp_{\mathrm{A}}\wedge\HypKosz\wedge\HypGen$, with
the class-$\mathsf M$ ambient strengthening above.
\end{example}

\subsection*{Containment partial order}

\begin{definition}[Package containment; \ClaimStatusDefinitional]
\label{def:type-containment}
For packages $\Hyp,\Hyp'$, write $\Hyp\leq\Hyp'$ if every primitive
appearing in $\Hyp$ appears in $\Hyp'$ (set inclusion of conjuncts),
or if $\Hyp'$ strengthens an ambient: $\HypAmb(\square)\leq
\HypAmb(\square')$ when $\square'$ refines $\square$ along the chain
$\Ch(\Vect)\leq\HypCoder\leq\HypWtComp\leq\Pro\leq\HypJad$. The
relation $\leq$ is a partial order: reflexive (any package contains
itself), antisymmetric (mutual containment forces equality of the
conjunct sets and ambient tags), transitive (transitive closure of
inclusion and ambient refinement).
\end{definition}

The ambient chain
$\Ch(\Vect)\leq\HypCoder\leq\HypWtComp\leq\Pro\leq\HypJad$ encodes the
class-$\mathsf M$ phenomenon: bar--cobar inversion in the raw
$\Ch(\Vect)$ ambient is genuinely false; passage to $\HypCoder$,
$\HypWtComp$, $\Pro$, or $\HypJad$ restores the equivalence. The
inscription cost rises with the refinement; using the weakest ambient
that closes the proof is the discipline (shadow-object distinction3).

\section{Composition along reconstruction theorems}
\label{sec:type-composition}

A reconstruction theorem at level interface $\BeilLevel\to\BeilLevel'$
bears a name $\Recon^{\BeilLevel\to\BeilLevel'}$ and a hypothesis
package $\Hyp_{\Recon}$. The Vol~I primary list is
\[
\begin{array}{l|c|l}
\textbf{Reconstruction} & \textbf{Levels} & \textbf{Package} \\
\hline
\text{Morita} &
 1\leftrightarrow 0 &
 \HypAug\wedge\HypMC \\
\text{Theorem A} &
 1\leftrightarrow 2 &
 \HypAug\wedge\HypComp\wedge\HypFinDimGr \\
\text{Theorem H} &
 1\to 3 &
 \Hyp_{\mathrm{A}}\wedge\HypKosz\wedge\HypGen \\
\text{Theorem B} &
 1\to 3 &
 \HypKosz\wedge\HypAmb(\HypCoder\joinp\HypWtComp\joinp\Pro\joinp\HypJad) \\
\text{Drinfeld double} &
 4\to 3 &
 \HypHall\wedge\HypIntegral\wedge\HypDescent \\
\text{Modular reconstruction} &
 5\to 4 &
 \HypMod\wedge\HypClut.
\end{array}
\]

\begin{proposition}[Composition rule]
\label{prop:type-composition}
\ClaimStatusProvedHere
If $\TypeSig_1=(\Quad_1,\Pres_1,\BeilLevel\to\BeilLevel',\Hyp_1)$
and $\TypeSig_2=(\Quad_2,\Pres_2,\BeilLevel'\to\BeilLevel'',\Hyp_2)$
are two reconstruction theorems whose target / source levels match,
their composition has type signature
\[
 \TypeSig_2\circ\TypeSig_1
 =\bigl(\Quad_1\meet\Quad_2,\ \Pres_1\meet\Pres_2,\ \BeilLevel\to\BeilLevel'',\ \Hyp_1\joinp\Hyp_2\bigr),
\]
where $\joinp$ is the join in the package partial order
(Definition~\ref{def:type-containment}) and $\meet$ on quadrant /
presentation tags is the intersection (refinement) of grid cells.
The hypothesis package of a composite reconstruction is the join of
the constituent packages; suppressing a conjunct of $\Hyp_1\joinp\Hyp_2$
in the inscription instantiates shadow-object distinction.
\end{proposition}

\begin{proof}
A theorem at level $\BeilLevel\to\BeilLevel'$ produces witnesses in the
target carrier under $\Hyp_1$; chaining a theorem
$\BeilLevel'\to\BeilLevel''$ requires its source hypotheses $\Hyp_2$ on
those witnesses. Both must hold, so the composite hypothesis is the
conjunction of the two; the partial order on packages identifies this
conjunction with the join. The quadrant / presentation tags refine to
the cells common to both reconstructions: a $\Open$-only theorem
composed with a $\mathrm{CY}$-only theorem has empty intersection and
no composition is defined.
\end{proof}

\begin{proposition}[Meet and join of packages]
\label{prop:type-meetjoin}
\ClaimStatusProvedHere
The package partial order
$\bigl(\{\Hyp\},\leq\bigr)$ admits binary meets and joins:
\[
 \Hyp\meet\Hyp'
 =\{\text{conjuncts in both}\},\qquad
 \Hyp\joinp\Hyp'
 =\Hyp\wedge\Hyp',
\]
extended on ambient tags by intersection / refinement of the
$\Ch(\Vect)\leq\HypCoder\leq\HypWtComp\leq\Pro\leq\HypJad$ chain. The
poset is a bounded lattice with bottom the empty package
$\bigl(\HypAmb(\Ch(\Vect))\bigr)$ and top the full package
$\Hyp_{\mathrm{full}}$ obtained by conjoining all primitives at the
strongest ambient tag $\HypJad$.
\end{proposition}

\begin{proof}
Conjunction is associative and commutative on a finite alphabet, so
binary meets and joins exist. The ambient tag chain is totally
ordered, so its meet / join is well-defined. Boundedness follows from
the finiteness of the primitive alphabet.
\end{proof}

The lattice structure makes the type system a constructive grammar:
every theorem inscription is parsed by reducing its hypothesis to a
package in the lattice; reconstruction-theorem composition is package
join.

\section{The Hypothesis-Package Lattice}
\label{sec:type-lattice}

The reconstruction theorems form the nodes of a partially ordered
diagram; an edge $\Recon\to\Recon'$ records that the target package
of $\Recon$ contains (refines) the source package of $\Recon'$, so
that $\Recon'$ is composable after $\Recon$.

\begin{figure}[ht]
\centering
\begin{tikzpicture}[
 node distance=14mm,
 every node/.style={
   draw,
   rounded corners=2pt,
   inner sep=4pt,
   align=center,
   font=\small
 },
 arrow/.style={-{Latex[length=2mm]}, thick},
 hypothesis/.style={font=\scriptsize, draw=none, rounded corners=0pt}
]

\node (morita)
 {Morita\\ \scriptsize $1\!\leftrightarrow\!0$\\
  \scriptsize $\HypAug\wedge\HypMC$};

\node[right=20mm of morita] (thmA)
 {Theorem A\\ \scriptsize $1\!\leftrightarrow\!2$\\
  \scriptsize $\HypAug\wedge\HypComp\wedge\HypFinDimGr$};

\node[below=12mm of thmA] (thmH)
 {Theorem H\\ \scriptsize $1\!\to\!3$\\
  \scriptsize $\Hyp_{\mathrm A}\wedge\HypKosz\wedge\HypGen$};

\node[left=20mm of thmH] (thmB)
 {Theorem B\\ \scriptsize $1\!\to\!3$\\
  \scriptsize $\HypKosz\wedge\HypAmb(\geq\HypCoder)$};

\node[below=12mm of thmH] (drin)
 {Drinfeld double\\ \scriptsize $4\!\to\!3$\\
  \scriptsize $\HypHall\wedge\HypIntegral\wedge\HypDescent$};

\node[left=20mm of drin] (modr)
 {Modular reconstruction\\ \scriptsize $5\!\to\!4$\\
  \scriptsize $\HypMod\wedge\HypClut$};

\draw[arrow] (morita) -- node[hypothesis, midway, above]
 {$+\HypComp\wedge\HypFinDimGr$} (thmA);
\draw[arrow] (thmA) -- node[hypothesis, midway, right]
 {$+\HypKosz\wedge\HypGen$} (thmH);
\draw[arrow] (thmA.south west) -- node[hypothesis, midway, above left]
 {$+\HypKosz, \HypAmb$} (thmB.north);
\draw[arrow] (thmH) -- node[hypothesis, midway, right]
 {$+\HypHall\wedge\HypIntegral$} (drin);
\draw[arrow] (drin) -- node[hypothesis, midway, above]
 {$+\HypMod\wedge\HypClut$} (modr);
\draw[arrow] (thmB.south) -- node[hypothesis, midway, left]
 {$+\HypMod$} (modr.north);

\end{tikzpicture}
\caption{Finite subdiagram of the hypothesis-package lattice on the
Open Beilinson-tower reconstruction theorems. Nodes carry the level interface and the
package; arrows record package containment under
Definition~\ref{def:type-containment}, labelled by the primitives
added along the edge. Composition along a directed path takes the
join of the edge increments; reconstruction composition also requires
matching level interfaces
(Propositions~\ref{prop:type-composition}
and~\ref{prop:type-lattice-wellformed}).}
\label{fig:type-lattice}
\end{figure}

\begin{proposition}[Package-lattice well-formedness]
\label{prop:type-lattice-wellformed}
\ClaimStatusProvedHere
The displayed finite package diagram above is a finite subdiagram of
the bounded hypothesis-package lattice. Along every directed path
the package contribution is the join of the displayed edge
increments. A directed path in this package diagram gives a composite
reconstruction theorem only after the level interfaces are separately
composable in the sense of Proposition~\ref{prop:type-composition}.
Package containment alone is not a functor between Beilinson levels.
\end{proposition}

\begin{proof}
Each edge realises a strict refinement of hypothesis packages, so
the displayed diagram is acyclic. Finiteness of the primitive
alphabet gives a bounded lattice by Proposition~\ref{prop:type-meetjoin}.
The join of edge increments is therefore independent of parentheses.
The level interface is an additional datum: a theorem of type
$i\to j$ can be composed with one of type $j\to k$, and no package
containment relation can replace this equality of the middle level.
Thus the displayed graph controls hypotheses, while
Proposition~\ref{prop:type-composition} controls reconstruction
composition.
\end{proof}

The type system fixes the inscription contract for the Master
Reconstruction Theorem (Vol~I climax): a Vol~I theorem with type
signature $(\Open,\Pres,\BeilLevel,\Hyp)$ admits a $\mathrm{CY}$
companion at the same level under the vertical-equivalence cells
listed in $\KSDual$, with both inscriptions named in the master
concordance (\S\ref{sec:master-concordance-cross-volume}).
